% GENERATED, edit the French unit or the translation table
\mxcorrections{9}\label{C09-CORR}
\textbf{Teacher use.} The answers below concern the tasks not worked in the lesson. First have students state their method; give graded help before the solution. The durations are estimates to be calibrated in class.
\subsection*{Diagnostic and remediation}
\begin{enumerate}
\item $2^7=128$. \textit{Action:} Rebuild the product of seven factors.
\item $\dfrac{1}{9}$. \textit{Action:} Go back to the reciprocal power.
\item $x=2$. \textit{Action:} Isolate the unknown while preserving equivalence.
\item $2x+2$. \textit{Action:} Apply the sum of the derivatives.
\item $\num{0,8}$. \textit{Action:} Subtract the rate from 1.
\item $5\times2^n$. \textit{Action:} Count the n multiplications.
\end{enumerate}
\subsection*{Investigation activity}
The condition $f\prime=f$ motivates $\Delta y\approx h y$. With h=$\dfrac{1}{4}$, the values are 1; \num{1,25}; \num{1,5625}; \num{1,953125}. This construction approximates the function and does not replace its characterisation.
\subsection*{Immediate applications}
\begin{enumerate}
\item $y=e+e(x-1)=ex$.
\item $e^{2x+1}$; $x+2>2x$, so $x<2$.
\item $g\prime(t)=-6e^{-2t}<0$ on $\R$: strictly decreasing.
\item $u_n=5(e^{\num{0,3}})^n$; the ratio over two units is $e^{\num{0,6}}$.
\item Apply the inequality proved with the argument $2x$; equality only at x=0.
\end{enumerate}
\subsection*{Exercises: method and results}
\paragraph{Exercise 9.1 -- Transforming.} $e^2$; $e^{5x}$; $1+e^{-x}$. For the last expression, distribute then simplify the products.
\paragraph{Exercise 9.2 -- Equations and inequalities.} By strict increase: $2x-1=x+3$, so x=4. Then $-3x\geq x-2$, that is $x\leq \dfrac{1}{2}$.
\paragraph{Exercise 9.3 -- Product and maximum.} $f\prime=(1-x)e^{-x}$; increasing up to 1 then decreasing. Maximum $f(1)=\dfrac{1}{e}$. Values at the endpoints 0 and $4e^{-4}$ are smaller.
\paragraph{Exercise 9.4 -- A modelled population.} 200; factor $e^{\num{0,1}}$; rate $e^{\num{0,1}}-1\approx\pc{10,52}$. Resources and the carrying capacity of the environment are not taken into account; this is a model over a limited time. Both curves requested pass through (0;1); the first increases, the second decreases. They are symmetric about the vertical axis; see the figure in the section From discrete to continuous.
\paragraph{Exercise 9.5 -- Searching for a threshold.} Test $n=0,1,\ldots$ in turn with an iteration limit. $Q(4)\approx\num{10,8268}\geq10$, $Q(5)\approx\num{6,5668}<10$; decrease proved, so the first integer is 5.
\paragraph{Exercise 9.6 -- A symmetry.} $f(-x)=f(x)$: even. $f\prime=e^x-e^{-x}$ has the sign of x, so minimum 2 at 0. Another proof: $f(x)-2=(e^{\dfrac{x}{2}}-e^{-\dfrac{x}{2}})^2$.
\paragraph{Exercise 9.7 -- Uniqueness under a hypothesis.} $g\prime=f\prime e^{-x}-fe^{-x}=0$. On $\R$, g is constant, and g(0)=1. Thus $f(x)=e^x$. This proof of uniqueness uses an exponential already constructed; it does not prove existence.
\paragraph{Exercise 9.8 -- Delayed review.} $\dfrac{u_{n+1}}{u_n}=e^2$; first term e. Sum $e(1+e^2+e^4)=\dfrac{e(1-e^6)}{1-e^2}$. The direct sum is just as efficient here.
\paragraph{CQFD, exercise 18 p.~187, questions a, b and c.} a: $e^{3x}$. b: $e^2$, because $e^{x+2}$ is divided by $e^x$, which is never zero. c: $e^{x+1}$. Do not confuse the number e with the variable x.
\paragraph{CQFD, exercise 30 p.~187, questions a, b and c.} a: $x\leq0$. b: $2x>1$, so $x>\dfrac12$. c: $-3x\geq2x+1$, so $x\leq-\dfrac15$. The strict increase of the exponential preserves the direction; dividing by $-5$ reverses it in c.
\subsection*{MCQ: answers and reasons}
\textbf{MCQ 1.}
\begin{enumerate}
\item B. Initial condition.
\item A. Addition property.
\item D. Factor -2.
\item C. Strict increase and $e^0=1$.
\item B. Positivity of the exponential.
\end{enumerate}
\textbf{MCQ 2.}
\begin{enumerate}
\item C. Compare the exponents.
\item A. Ratio of consecutive terms.
\item D. Value 1 and slope 1.
\item B. Product rule.
\item C. Minimum of the difference.
\end{enumerate}
\subsection*{True or false}
\begin{enumerate}
\item False: at 0, 1 is different from 2.
\item True: its derivative $-3e^{-3t}$ is negative.
\item False: the discrete rate is $e^{\num{0,1}}-1$, about \pc{10,52}.
\end{enumerate}
\paragraph{Differentiated teaching.} For a technical block, offer only the first step (domain, factor or reference sample space), then withdraw the help. For an independent student, ask for a second method and the treatment of a limiting case. Orally, allow an initial reformulation in the language the student has mastered, then require a precise mathematical sentence in French.
